\SourceLeaf{379}%
in jeder Hinsicht vorsichtiger vor. Mit Pulversignalen auf dem Mont Tabor verband man den Mont Cenis mit dem Mont Colombier, mit Pulversignalen auf dem Berg Pierre sur autre den Mont Colombier mit dem französischen Dreieckspunkt Puy d’Usson. Zugleich gab man auch auf dem Mont Colombier selbst Signale, um diesen Ort mit der Genfer Sternwarte zu verbinden. Auf dem Mont Cenis und dem Mont-Colombier bestimmte man die Zeit aus Beobachtungen an Passagenfernrohren von \textsc{Lenoir}\footnote{Nach einigen Umständen zu schließen, scheint man 1821 und 1822 auf dem Mont Cenis dasselbe Passagenfernrohr verwendet zu haben, obwohl hier ein anderer Hersteller genannt wird.} und \textsc{Grindel}. Am französischen Standort gründete die Zeitbestimmung auf absoluten Sternhöhen, die mit einem Repetitionskreis gemessen wurden (vgl. \textit{Conn. des tems} 1829 und unsere Anzeige vom 10. Januar 1828). Auf dem Mont Cenis musste man sich jedes Mal drei Wegstunden vom Hospiz entfernen, um die Signale sehen zu können. Schließlich finden wir hier noch die Bestimmung des Längenunterschieds zwischen den Sternwarten von Turin und Mailand durch Pulversignale auf dem S. Bernardo di Fenera zu drei verschiedenen Zeiten \(1823\mathpunct{.}\mathpunct{.}\mathpunct{.}1824\). Dabei waren alle Umstände so günstig, wie sie bei Arbeiten dieser Art nur sein können.

Der dritte Abschnitt enthält die Breitenbestimmungen der Sternwarten auf dem Mont Cenis, auf dem Mont Colombier und in Turin, die beiden ersten mit Repetitionskreisen von \textsc{Troughton} und \textsc{Reichenbach}, die letzte mit dem Meridiankreis von \textsc{Reichenbach}. Die letzten Beobachtungen erreichen nicht ganz den Grad an Übereinstimmung, den man sonst von diesen Instrumenten gewohnt ist. Die Verfasser haben selbst darauf hingewiesen. und lassen offen, ob solche Anomalien wirklich bestehen oder auf einem Fehler bei der Handhabung des Instruments beruhen, der sich künftig aufklären lässt. Der Rezensent räumt ein, dass bei den vielen Vorsichtsmaßnahmen, die dieses Instrument verlangt, niemand, der nicht vor Ort ist, auch nur eine plausible Vermutung dazu äußern kann. Die Art, wie sich die Verfasser zu den Anomalien geäußert haben, veranlasst ihn jedoch, aus eigener Erfahrung ein Beispiel dafür anzuführen, wie geringfügige Umstände Beobachtungen manchmal beeinträchtigen können. Im Laufe mehrerer Jahre war die schöne Übereinstimmung der Beobachtungen an einem Meridiankreis, der dem Turiner vollkommen glich, nur ein einziges Mal
